\begin{abstract}
Large EEG foundation models combine domain-specific architectures with
thousands of hours of self-supervised pretraining.  We ask whether
off-the-shelf vision pretraining can provide a ``free lunch'' for EEG
decoding without EEG-specific pretraining.  The answer depends on the geometry
presented to the visual encoder.  We introduce \method{}, a parameter-free,
bijective adapter that folds a raw EEG matrix from $C\times T$ to
$(CP)\times(T/P)$.  Folding preserves every sample while placing consecutive
temporal phases on adjacent rows, exposing structured multi-lag neighborhoods
and retaining more signal-bearing height through hierarchical downsampling.
We fine-tune a one-channel EfficientNet-B0 end to end, without spectral
conversion or EEG pretraining, on 11 benchmarks spanning clinical, affective,
sleep, and BCI applications. Under validation checkpoint selection and
repeated-seed evaluation, we report balanced accuracy, PR-AUC/ROC-AUC, or
$\kappa$/weighted F1 in one metric protocol alongside BIOT, LaBraM, \cbramod{},
and \reve{} references. The compact visual model is competitive on several
BCI and clinical tasks but remains below EEG foundation models on
topology- and subject-shift-sensitive cases such as FACED. These results do
not establish universal superiority or replace same-pipeline reproductions.
They show that representation geometry is a necessary control when attributing
downstream gains to EEG foundation-model pretraining. Matched five-seed controls on TUEV
and PhysioNet-MI further show that visual pretraining improves all three tested
backbones over random initialization, providing direct evidence of visual
pretraining transfer to EEG under end-to-end fine-tuning.
\end{abstract}
